%%%%%%%%%%%%%%%%%%%%
%%%% FORMATTING %%%%
%%%%%%%%%%%%%%%%%%%%
\documentclass[11pt,a4paper]{article}  % Set the class of the document to article with 10pt font size and A4 paper size
\setlength{\textheight}{25.7cm}  % Set the height of the text area to 25.7cm
\setlength{\textwidth}{16cm}  % Set the width of the text area to 16cm
\setlength{\unitlength}{1mm}  % Set the default unit length for picture environments to 1mm
\setlength{\topskip}{2truecm}  % Set the skip above the top baseline of the text area to 2cm
% Calculate and set the top margin
\topmargin 260mm \advance \topmargin -\textheight 
\divide \topmargin by 2 \advance \topmargin -1in 
% Set the header height and separation to 0pt
\headheight 0pt
\headsep 0pt
% Calculate and set the left margin
\leftmargin 210mm \advance \leftmargin -\textwidth 
\divide \leftmargin by 2 \advance \leftmargin -1in 
% Set the odd and even side margins to the calculated left margin
\oddsidemargin \leftmargin
\evensidemargin \leftmargin
% Set the paragraph indentation to 0pt
\parindent=0pt
% Formatting of (sub)section headers
\usepackage[medium]{titlesec}
\titlespacing\section{0pt}{12pt plus 4pt minus 2pt}{0pt plus 2pt minus 2pt}
\titlespacing\subsection{0pt}{12pt plus 4pt minus 2pt}{0pt plus 2pt minus 2pt}
\titlespacing\subsubsection{0pt}{12pt plus 4pt minus 2pt}{0pt plus 2pt minus 2pt}
% Paragraph formatting
\setlength{\parindent}{0ex}  % Removes default indent at new paragraph
\setlength{\parskip}{.5em}  % Sets the vertical space between two paragraphs.
% Move down page numbers
\setlength{\footskip}{2cm} 
% Include subsubsections in table of contents
\setcounter{tocdepth}{3}


%%%%%%%%%%%%%%%%%%
%%%% PACKAGES %%%%
%%%%%%%%%%%%%%%%%%
\usepackage{tikz}
\usepackage[toc, page]{appendix}  % For adding appendices.
\usepackage[english]{babel}  % Provides multilingual support for LaTeX documents.
\usepackage[format=plain, font=it]{caption}  % Customizes the formatting of captions for figures and tables.
\usepackage[utf8]{inputenc}  % Specifies the input encoding of the document to UTF-8, allowing for the direct input of special characters.
\usepackage{amssymb}  % Provides access to various mathematical symbols and fonts.
% \usepackage{commath}  % Offers additional mathematical commands and environments for typesetting mathematical expressions.
\usepackage{mathtools}  % Enhances the functionality of mathematical typesetting, providing additional tools and symbols for mathematical equations.
\usepackage[shortlabels]{enumitem}  % Customizes the appearance and behavior of lists, allowing for different label styles and formatting options.
\usepackage{gensymb}  % Provides access to generic symbols, including common symbols like degree symbol (°).
\usepackage{calligra}  % Allows the use of the calligraphic font style in the document.
\usepackage{amsmath}  % Enhances the functionality of mathematical typesetting in LaTeX.
\usepackage[hidelinks]{hyperref}  % Enables the creation of hyperlinks within the document, with the hidelinks option used to hide the colored boxes around hyperlinks.
\usepackage{wrapfig}  % Facilitates the wrapping of text around figures or tables.
\usepackage{booktabs}  % Enhances the quality of tables by providing guidelines for proper table formatting.
\usepackage{multicol}  % Enables the creation of multi-column layouts in the document.
\usepackage{multirow}  % Allows for the merging of rows in tables.
\usepackage{ctable}  % Provides customizable tables with improved features compared to standard LaTeX tables.
\usepackage{subcaption}  % Enhances the functionality of subcaptions for subfigures and subtables.
\usepackage{systeme}  % Aids in the typesetting of systems of equations.
\usepackage{lipsum}  % Generates dummy text (lorem ipsum) for testing the layout and formatting of the document.
\usepackage{tocloft}  % Allows for the customization of table of contents formatting.
\usepackage{secdot}  % Adds a dot after section numbers in the document.
\usepackage{dirtytalk}  % Typeset quotations with proper quotation marks
\usepackage{systeme}  % For systems of equations 
\usepackage{esint}  % Nice looking integrals
\usepackage{enumitem}  % Offers more flexibility for lists
\usepackage{wrapfig}  % Wrap text around figures


% Style for code using listings
\usepackage{listings}
\usepackage{xcolor}

\definecolor{codegreen}{rgb}{0,0.6,0}
\definecolor{codegray}{rgb}{0.5,0.5,0.5}
\definecolor{codepurple}{rgb}{0.58,0,0.82}
\definecolor{backcolour}{rgb}{0.95,0.95,0.92}
\lstdefinestyle{mystyle}{
	backgroundcolor=\color{backcolour},   
	commentstyle=\color{codegreen},
	keywordstyle=\color{magenta},
	numberstyle=\tiny\color{codegray},
	stringstyle=\color{codepurple},
	basicstyle=\ttfamily\footnotesize,
	breakatwhitespace=false,         
	breaklines=true,                 
	captionpos=b,                    
	keepspaces=true,                 
	numbers=left,                    
	numbersep=5pt,                  
	showspaces=false,                
	showstringspaces=false,
	showtabs=false,                  
	tabsize=2
}
\lstset{style=mystyle}


% Set colors for citation numbers, equation numbers, links and URLs
% Blue citation numbers
% Magenta links (e.g. in table of contents)
% Red equation numbers
\hypersetup{colorlinks, linkcolor={black}, citecolor={blue}, urlcolor={magenta}} 
\renewcommand\theequation{{\color{red}\arabic{equation}}} 

% Symbols for footnotes
\usepackage[symbol]{footmisc}
\renewcommand{\thefootnote}{\fnsymbol{footnote}}

% Table and Figure captions formatting
\usepackage[font={sl}, format = hang, labelfont=bf]{caption}  


%%%%%%%%%%%%%%%%%%
%%%% COMMANDS %%%%
%%%%%%%%%%%%%%%%%%
\renewcommand{\d}[1]{\ensuremath{\operatorname{d}\!{#1}}} % \d command for upright d in differentials. Also has some space before it for at the end of an integral.
\newcommand{\angstrom}{\text{\normalfont\AA}}  % \angstrom argument for angstrom symbol
% Below is for recreating the script-r from Griffith's "Introduction to Electrodynamics".
% Use \scripty{r} as command. Can also be made bold or given a hat.
\DeclareMathAlphabet{\mathcalligra}{T1}{calligra}{m}{n}
\DeclareFontShape{T1}{calligra}{m}{n}{<->s*[2.2]callig15}{}
\newcommand{\scripty}[1]{\ensuremath{\mathcalligra{#1}}}  


%%% Color boxex for "Optional" and "Example" parts.
\usepackage{tcolorbox}
\tcbuselibrary{skins, breakable}
\definecolor{darkblue}{RGB}{0, 0, 139}
\definecolor{darkgreen}{RGB}{0, 80, 0}
\definecolor{warningorange}{HTML}{B06000}

\newtcolorbox{optionalbox}[1][]{
	breakable,
	enhanced,
	colback=gray!8,
	colframe=gray!50,
	colupper=darkblue,
	fontupper=\normalfont,
	fonttitle=\bfseries\color{darkblue},
	title=Optional,
	attach boxed title to top left={yshift=-2mm, xshift=4mm},
	boxed title style={colback=white, colframe=gray!50},
	#1
}

\newtcolorbox{examplebox}[1][]{
	breakable,
	enhanced,
	colback=gray!8,
	colframe=gray!50,
	colupper=darkgreen,
	fontupper=\normalfont,
	fonttitle=\bfseries\color{darkgreen},
	title=Example,
	attach boxed title to top left={yshift=-2mm, xshift=4mm},
	boxed title style={colback=white, colframe=gray!50},
	#1
}


\usepackage{menukeys}

\usepackage{graphicx}
\usepackage{array}
\usepackage{longtable}
\usepackage{makecell}
